\section{Experimental Design}
\label{sec:experimental-design}
The objective of the evaluation is to determine when the authenticated
cost-change mechanism reduces proof-service work while retaining a verified
decision about the original quantized objective. We examine the contribution
of its components, the cost of reuse, and the workloads for which the compact
interval remains useful. This version reports completed local measurements.
The additional server study is prospective and contributes no result to
the tables below.

\subsection{Research Questions}
\label{sec:ccd-research-questions}
We organize the evaluation around four research questions:
\begin{itemize}[leftmargin=*]
\item \textbf{RQ1: Component contribution.}
How do exact old groups, switching information, the fitted basis and bounded
row refinement affect decision resolution and representation quality?
\item \textbf{RQ2: Reuse and total cost.}
What are the separate preparation and request costs, and when does sharing
authenticated state reduce measured work relative to the implemented controls?
\item \textbf{RQ3: Computational scale.}
How do catalog size, authentication, proof generation, public verification
and proof material affect the observed computational cost?
\item \textbf{RQ4: Operating limits.}
How do difficult candidates and reference changes affect reuse, and can
cost-aware scheduling avoid unsuccessful certificate work?
\end{itemize}

\subsection{Datasets}
\label{sec:ccd-datasets}
The earlier seed-521 pilots use Fashion-MNIST images~\cite{fashion2017} and frozen
384-dimensional neural text features from 20 Newsgroups
\cite{news20data,minilmcard2026}. Images have $d=784$ and both workloads use
$k=5$ centers. The text extractor is the archived
\texttt{all-MiniLM-L6-v2} model, with attention-mask mean pooling and
normalization. Its weights, selected document indices and embedding files
are pinned. The proof statement starts from the resulting integer features;
it does not prove encoder inference.

Each pilot fits the reference and representation on a 512-record past-training
partition. A separate 512-record incoming-training partition produces candidate
banks, and a held-out feature pool supplies the audited rows. The deployed
reference is the eighth training Lloyd iterate. The fitted configuration uses
seed 521. We retain catalog prefixes of $n=32$ and $n=128$ for real proofs;
native checks also use the retained larger pools. These sources and fitted
workloads have been inspected during development and are therefore treated
as exploratory evidence, not independent test-set confirmation.

The later published-baseline gate independently refits each of Fashion,
text and UCI HAR ($d=561$) with seeds 557, 563 and 569. Each fit uses
512 past-training and 512 incoming-training records, with catalog prefixes
of 32 and 128 held-out rows. All first eight proposals from the first two
incoming minibatches and the four fixed update rates are retained. Native
quality therefore contains 24 requests per source/catalog size across
three fits. Real proofs retain all first four requests from seed 557 on
Fashion and text with $n=32$. These additional fits measure configuration
variation on inspected sources, not untouched external confirmation.

The manuscript also retains completed earlier controls: a 1,024-record
Fashion workload, repeated MNIST timing controls~\cite{lecun1998}, and native
checks on UCI HAR sensor features~\cite{ucihar2012}. They support the
previously stated mechanism and application observations. Their fitted
configurations and execution cohorts are identified separately; they are
not pooled with the new pilots as additional independent repetitions.

\subsection{Baselines}
\label{sec:ccd-baselines}
\paragraph{Exact Full}
The internal control proves exact original-space nearest-center costs using
Groth16 over BN254~\cite{groth2016}. It receives the same catalog integers,
reference and candidate banks. Its public-center circuit and proved reference
cost can be reused. Since an exact public cost reveals more than an interval
or a sign, this control does not establish optimality against every possible
exact sign-only implementation.

\paragraph{PCA certificate}
PCA replaces the centroid-informed fitted basis while retaining the same
authentication, switching relation, correction cap and exact fallback.
The comparison isolates the representation within our certificate framework.
The wrapper is not presented as a separate published PCA proof system.
Policy comparisons require the same scheduling optimization on both bases.

\paragraph{ZKML primitive adapter}
The external-backend control adapts the exact dot-product, adder and Poseidon
primitives from ZKML~\cite{chen2024zkml}, pinned at upstream commit
\texttt{50d0cf30e2bbcb8b1f31839b775c1d018b720527}. The adapter implements
the same integer nearest-center objective with Halo2--KZG and an
owner-approved commitment to the same rows. It is not the unmodified
optimizing compiler. Its commitment and public exact-cost output differ
from our Groth16 interval interface, so backend comparisons concern the
implemented systems rather than an isolated projection benefit.

\paragraph{Paper-derived numerical controls}
We implement the random-sign embedding of Boutsidis et
al.~\cite{boutsidis2010random} and the sparse three-point embedding of
Achlioptas~\cite{achlioptas2003database}, each at empirical prediction rank
63. Their supplied-bank costs use raw projected distances, without refitting
clusters. The lightweight coreset of Bachem et
al.~\cite{bachem2018scalable} uses 16 or 32 draws with replacement and
inverse-probability weights; uniform sampling is an additional control.
These budgets do not assert particular theorem-level approximation
parameters or supply an authenticated full-catalog sign certificate.

A separate follow-up implements the one-bit RaBitQ estimator of Gao and
Long~\cite{gao2024rabitq}, with exact old groups as IVF anchors and an
exhaustive scan of all row/center pairs. It is a float64 reference, without
the published ANN engine's SIMD or query-quantization optimizations. A
spherical-coordinate quantile envelope assigns a union-bound failure
budget of 0.01 to the predeclared finite pair family per fit/catalog case.
The rotation is independent of the fixed candidates. This is an
ideal-arithmetic statistical interval, without outward-rounding proof or
cryptographic authentication. Exact old cost is supplied by numerical
old-group preprocessing.

Separate sign/sparse projection certificate arms attach our existing
integer relation and correction policy. Their residual-aware predictor
differs from raw projected-distance prediction. They are adaptations,
not published proof systems; weak coverage of this particular envelope
does not refute the original randomized approximation results.

\paragraph{Component controls}
Authenticated moments-only and advance-bank affine controls are retained
where complete measurements already exist. The former omits row-level
switching information. The latter knows the union of future centers before
preprocessing and therefore has a different information contract. Numerical
rank and correction-budget studies are explicitly distinguished from
proof-generation measurements.

\subsection{Experimental Protocol}
\label{sec:ccd-protocol}
The catalog $X$, deployed bank $C^0$, basis and quantizer remain fixed within
one reference episode. Candidates $C^1,C^2,\ldots$ can change every center.
Sixteen chronological incoming minibatches of 32 records, combined with
update rates $0.125$, $0.25$, $0.5$ and $1$, define 64 proposals in batch-major
order. The eight-request pilot keeps the first eight; the text follow-up
keeps the first sixteen. Candidate construction does not use audit outcomes.
Exact public duplicates may share work only under the same rule for every
method, and total and distinct request counts remain recorded.

Every conclusive interval is checked against an independently recomputed
integer cost change. A valid interval spanning zero remains inconclusive;
after at most 16 proved row corrections, unresolved requests generate and
verify new Full proofs. Failed certificate work is retained in the total.
The auditor replays public-only bundles against separately retained model,
root, challenge, task and verification-key pins.

We distinguish two execution settings. A \emph{cold campaign} total
includes the recorded preparation, key construction, initial evidence,
request proofs, replay and fallback. Recorded fitting is retained where
charged by the source campaign; the text16 extension supplies frozen fitted
representations. A \emph{warm} trial reuses the applicable setup, authenticated
rows and reference-cost evidence for both arms, but creates new request
proofs. Copying/export and independent numerical validation are outside the
warm phase-wall sums. These phase totals are not independently timed RPC
latencies.

\subsection{Experimental Settings}
\label{sec:ccd-settings}
The new measured comparisons run serially in Linux/WSL on the same laptop,
with eight proof workers and CPU affinity fixed to eight logical processors.
The process-tree RSS cap is 3.5~GiB. A separate Windows host-memory guard
can stop a run when the operating system signals pressure; such a stop is
retained as an incomplete measurement. Native BLAS is restricted to one
thread during these campaigns. No GPU is used for the audited proof workload.

The default prediction rank is $m=63$, certificate rank is $s=79$, and basis
scale is $T=2^{14}$. Coordinates use the admitted 12-bit integer profile.
Compact switching batches contain 16 rows, while Full query batches contain
eight. The correction cap is also 16, independently of the switching batch
size. The new pilots have one timing execution per fitted configuration;
the repeated historical MNIST controls belong to their own archived cohorts.

The cost-aware follow-up estimates proof work from the earlier eight-request
pilot. It plans width-ordered corrections using the prover's private rows
before producing query proofs. A batch activates compact proofs only when
the predicted saving also covers key loading and a ten-percent margin.
These native values select witness work; they do not certify the answer.
The public verifier still requires disjoint, complete request coverage and
valid selected proofs. This scheduler is an engineering variant of the
method in Section~\ref{sec:methodology}, not a new certificate relation.

In the later gate, every representation receives the same final routing
policy and the same seed-521 calibration, frozen before its own warm
measurements. The four-request batch pays its full native preflight.
All changing-request proofs are fresh, and one newly timed Full arm is
shared per dataset. The study has one proof timing execution per cell;
three within-fit native timing passes are not three independent fits.
The main gate was frozen before its outcomes. RaBitQ and the warm-routing
follow-up were registered after the first native outcomes and before
their respective outcomes, so the combined study remains exploratory.

\subsection{Evaluation Metrics}
\label{sec:ccd-metrics}
We report three separate quality measures. \emph{Decision coverage} is the
fraction of requests with a conclusive base or refined interval before Full.
\emph{Assignment agreement} compares approximate nearest-center assignments
with exact assignments on the same quantized rows. \emph{Label-certificate
coverage} counts individual assignments justified by their pairwise bounds.
None of these is ground-truth class-label accuracy. Final service decisions
must agree with the exact cost ordering after all required verified work.

For numerical cost estimators, estimated-sign agreement compares the
estimated and exact changes. Balanced sign agreement averages recall over
the observed true sign classes and is reported descriptively to expose
class imbalance. It was added after inspecting the frozen outcomes.
Coresets have no label output for unsampled rows, so their label agreement
and deterministic certificate coverage are not applicable. Statistical
interval resolution is reported separately from deterministic coverage.

For efficiency, we report phase-wall and enclosing totals according to their
recorded scope, the number of corrections and Full requests, peak RSS, and
proof bytes. Fresh request proof bytes are separated from complete exported
proof bytes containing reusable row evidence. Keys, parameters and public
metadata are not silently counted as zero storage.

To expose amortization, we further separate the preparation term in
Eq.~\eqref{eq:methodology-cost}. Let $K_{\mathrm{setup}}$ denote circuit/key
construction, $P_{\mathrm{epoch}}$ the remaining epoch preparation,
$L(Q)$ recurring process/key loading, and $R_t$
the complete request work, including any corrections and exact evaluation:
\begin{equation}
 T(Q)=K_{\mathrm{setup}}+P_{\mathrm{epoch}}+L(Q)+\sum_{t=1}^{Q}R_t.
 \label{eq:ccd-accounting}
\end{equation}
The per-request moment-link proof is charged in $R_t$ in the current
implementation. The empirical crossing with baseline $b$ is the first
measured prefix satisfying $T_{\mathrm{cert}}(Q)<T_b(Q)$. A crossing derived
from constant mean costs is a model estimate, not a measured threshold.

